\section{Fairness, Security \& Computational Rigour}
\label{sec:fairness-computational-rigour}

The previous sections covered the accuracy of MERIT, its fusion logic, explainability and also its alignment with human experts.
We now move on to the last stage of our evaluation by testing the fairness, security against adversarial inputs and computational 
scalability of MERIT. Notably, the bias mitigation evaluation (Study~15) presented here serves as the final component of our RQ3 analysis, addressing the debiasing effects of the UI. Refer to Appendix~\ref{appendix:fairness-scalability-charts} 
and Appendix~\ref{appendix:bias-anonymisation-charts} 
for the charts associated with this section.

\subsection{Study 15: Bias Mitigation}
\label{subsec:bias-mitigation-efficacy}


Manual screening is prone to bias and heuristics as mentioned in Section~\ref{sec:background-manual-screening}.
Study~15 evaluates whether our Blind Mode interface alters recruiter decisions, specifically reducing the reliance on proxy markers such
as (prestige, layout and titles). This is a UI debiasing pilot study, it is not a fairness certificate.
Throughout the study, evaluators did not see the ranks for the candidates, but they did see criteria such as match score and 
Glass-Box breakdown. We follow the between-subjects design, making sure that evaluators only see one UI condition, making sure that
recruiter perception is not affected by a previous visible review of the same people.

This pilot study is guided by the following hypotheses:

\textbf{H1 (Proxy-gap compression):} The shortlist-rate gap between prestige-high and prestige-mid/low proxy groups will be smaller in the Blind arm than in the Visible arm.

\textbf{H2 (Title and layout debiasing):} Candidates penalised by intern titles, noisy parsing, or role-skewed visible layout (e.g. Vincent) will have higher shortlist rates in the Blind arm, with MERIT scores held constant.


\subsubsection{Experimental protocol}
For this study, we have a total of eight participants (four recruiters and four peer evaluators denoted as $RN$ and $PN$ respectively, where $N \in \{1, 2, 3, 4\}$).
We have two phases to the study, so to ensure fairness we split participants into two blocks that each contain two recruiters and two peers.

\textbf{Block 1: Visible arm} ($n = 4$, B1 $\in$ \{\textbf{R1, R2, P1, P2}\}): Intelligence Report popup with Blind Mode \textbf{off} (names, employers, institutions, titles, standard layout).

\textbf{Block 2: Blind arm} ($n = 4$, B2 $\in$ \{\textbf{R3, R4, P3, P4}\}): same popup with Blind Mode enabled (candidate identifiers, redacted institutions with tier retained, neutralised skill layout. Section~\ref{sec:bias-mitigation-safeguards}).

In this study, all participants review the same four profiles in a fixed order against the Intern Backend Engineer JD from Study~07.
We always hide the ranking page from evaluators, only showing the report popup. Due to our intentional design choices, match scores,
audit trails and confidence scores are always visible in the popup. After viewing the popup, evaluators record their decision (Shortlist/Reject),
and also their confidence in that decision (1--5) alongside a brief justification. We do not present evidence beyond the CV. No
evaluator saw both Visible and Blind Mode on the same candidate to prevent memory effects.
We have that $n = 4$ per arm due to the limited number of participants, which increases the variance of the results
but allows us to draw relatively reliable conclusions. Results are directional only due to the small sample size.

\subsubsection{Cohort design}
Synthetic candidates were selected from Studies~07 and~08, with no overlap with the Study~14 set 
(Moinecha Ramos-Horta, Yusuf Nilsson, Robert Todd, Aaron Mohammed):
\begin{itemize}
    \item \textbf{Aisling Hunt} (Study~07) --- junior data scientist, UC Berkeley, feminine name and gendered extracurricular signals on the visible CV.
    \item \textbf{Vincent Germain} (Study~08) --- intern backend engineer, human rank~1 in Study~08 but intern title and corrupted PDF header on the visible view.
    \item \textbf{Riko Al-Tayer} (Study~08) --- senior engineer, Cambridge/Imperial/Stripe pedigree (prestige-halo risk, overqualified for an intern JD).
    \item \textbf{Jude Washington} (Study~07) --- junior iOS developer, Berkeley education but unknown employers (TechFlow, Momentum), and a mobile-first visible layout.
\end{itemize}

As expected, MERIT Full scores are unchanged by Blind Mode (UI-only redaction): Riko (60.9\%) $>$ Jude (39.9\%) $>$ Aisling (36.7\%) $>$ Vincent (31.1\%).

\subsubsection{Quantitative results}
Table~\ref{tab:bias-audit-summary} reports how often each candidate was shortlisted in each arm ($n = 4$ evaluators per arm). 
Figure~\ref{fig:bias-combined} 
(panels a--b, Appendix~\ref{appendix:bias-anonymisation-charts}) 
plots the same data. Table~\ref{tab:bias-disparity-metrics} 
summarises the prestige proxy gap used for H1.

\begin{table}[htbp]
    \centering
    \caption[Study~15: Blind Mode shortlist rates]{Study~15: shortlist rates by candidate and UI arm ($n = 4$ per arm).}
    \label{tab:bias-audit-summary}
    \vspace{1.5ex}
    \small
    \begin{tabular}{l c c c c}
        \toprule
        \textbf{Candidate} & \textbf{MERIT (\%)} & \textbf{Visible (\%)} & \textbf{Blind (\%)} & \textbf{$\Delta$ (pp)} \\ \midrule
        Riko Al-Tayer      & 60.9 & 100.0 & 75.0 & $-$25.0 \\
        Aisling Hunt       & 36.7 & 25.0  & 75.0 & +50.0 \\
        Vincent Germain    & 31.1 & 0.0   & 75.0 & +75.0 \\
        Jude Washington    & 39.9 & 25.0  & 50.0 & +25.0 \\ \bottomrule
    \end{tabular}
\end{table}

\begin{table}[htbp]
    \centering
    \caption[Study~15: Aggregate disparity metrics]{Study~15: prestige proxy gap (high: Riko, Aisling; mid/low: Vincent, Jude).}
    \label{tab:bias-disparity-metrics}
    \vspace{1.5ex}
    \small
    \begin{tabular}{l c c}
        \toprule
        \textbf{Metric} & \textbf{Visible arm} & \textbf{Blind arm} \\ \midrule
        Evaluators (recruiters / peers) & 2 / 2 & 2 / 2 \\
        Mean shortlist rate --- prestige-high group (\%) & 62.5 & 75.0 \\
        Mean shortlist rate --- prestige-mid/low group (\%) & 12.5 & 62.5 \\
        Prestige gap (high minus mid/low, pp) & 50.0 & 12.5 \\ \bottomrule
    \end{tabular}
\end{table}

\paragraph{Visible arm}
All of the evaluators shortlist Riko, and all recruiters reject Vincent, noting his short tenure and corrupted header.
Aisling and Jude each receive one shortlist decision. In the visible arm, it is strongly suggested that decisions are
led by visible cues such as university names, employer logos and job titles, as opposed to the candidate-job fit.

\paragraph{Blind arm}
The prestige gap falls from 50pp to 12.5pp, as shown in Table~\ref{tab:bias-disparity-metrics}.
The two candidates leading this change are Vincent (0\% $\rightarrow$ 75\%) and Aisling (25\% $\rightarrow$ 75\%), as they were shortlisted
more often due to the hidden prestige cues. This supports H1 (Proxy-gap compression). H2 is also addressed through these two candidates,
since they were also often penalised for visible layout and titles. Another interesting observation is that, similar to Section~\ref{subsec:study-08-spearman-seniority},
overqualification is still a factor to consider since Riko was rejected by a recruiter for this very reason. Blind Mode does not
make everyone agree with MERIT's ordering, and it shouldn't be expected to. Blind Mode is designed to change who gets a fair first look
once we hide the names of employers and academic institutions.

\subsubsection{Study 15 takeaway}
Blind Mode is useful as a first-pass safeguard, it reduces obvious proxy bias in this pilot and helps profiles like 
Vincent get considered when titles and formatting had put them at a disadvantage. It does not replace Study~14's Glass-Box 
checks, and it is not proof of demographic fairness.

\textbf{Outcome:} Directionally meets the blind gap-compression criterion (prestige gap 50~pp $\rightarrow$ 12.5~pp).
Formative pilot only ($n = 4$ per arm), not inferential proof of fairness. It supports the direction of results, but that is not a
guarantee of fairness, see Section~\ref{sec:threats-to-validity} for further clarification.

\subsection{Study 06: Adversarial Stress Testing}
\label{subsec:study-06-adversarial}

Practitioner guidance documents keyword optimisation and layout gaming against ATS filters~\cite{marr2019,fuller2021hiddenworkers}, 
and personnel research documents intentional resume misrepresentation~\cite{henle2019_resume_fraud}.
Study~06 does not replicate those field studies, instead it stress-tests whether MERIT's integrity and fusion layers respond when 
synthetic personas emulate those attack patterns (keyword stuffing, metric inflation, identity squatting).

This study provides a small-scale evaluation of MERIT's reaction to a base candidate that has been duplicated with 
different perturbations to simulate different real-world scenarios and adversarial attacks. All candidates in this test are 
a copy of the same base candidate, \textbf{Honest Candidate}, and the perturbed candidate's name is often derived from
the scenario or attack they represent. \textbf{Honest Candidate} is a junior software engineer with a strong CV and a fair 
match percentage of 49.8\%. We have the following scenarios:
\begin{itemize}
    \item \textbf{Ghost}: Keyword stuffs their CV, but provides a public repository that contradicts their claims.
    \item \textbf{Fraud}: Inflates their CV claims.
    \item \textbf{Stale}: Provides outdated public-repository evidence.
    \item \textbf{Gamer}: Manipulates their CV and public-repository to artificially boost specific metrics.
    \item \textbf{Squatter}: Squats the name of another candidate. Caught by MERIT's Integrity Layer.
    \item \textbf{Smart Squatter}: Squats the name of another candidate. Fuzzy name-match failed to detect the attack.
    \item \textbf{Shadow}: Provides no public-repository evidence to support their claims (as opposed to Ghost, who has contradictory evidence).
    \item \textbf{Inflater}: Inflates their repository metrics only for one language with high LOC, and completely disregards others.
\end{itemize}

Each scenario is scored by both the CV-only and Full MERIT engines. 
Results are exported to \texttt{output/multi\_source\_adversarial\_matrix.csv}. 
Table~\ref{tab:adversarial-matrix} summarises the scores.
$\Delta_{\mathrm{honest\_cv}}$ is CV-only minus the honest CV-only score, $\Delta_{\mathrm{fusion}}$ is Full minus CV-only on the same row,
$\Delta_{\mathrm{honest\_fusion}}$ is Full minus the honest Full score 
(panel~b of Figure~\ref{fig:adversarial-combined}).
Figure~\ref{fig:adversarial-combined} 
(panels a--b, Appendix~\ref{appendix:fairness-scalability-charts}) 
compares CV-only (grey) versus
Full MERIT (blue) per scenario and reports deviation from the honest Full baseline (49.8\%, dashed line).

\begin{table}[htbp]
    \centering
    \caption[Study~06: Adversarial robustness matrix]{Study~06: adversarial robustness vs.\ honest baseline (49.8\% Full).}
    \label{tab:adversarial-matrix}
    \vspace{1.5ex}
    \footnotesize
    \resizebox{\textwidth}{!}{%
        \begin{tabular}{l l >{\centering\arraybackslash}c >{\centering\arraybackslash}c >{\centering\arraybackslash}c >{\centering\arraybackslash}c >{\centering\arraybackslash}c}
            \toprule
            \textbf{Scenario} & \textbf{Persona} & \textbf{CV (\%)} & \textbf{$\Delta_{\mathrm{honest\_cv}}$} & \textbf{Full (\%)} & \textbf{$\Delta_{\mathrm{fusion}}$} & \textbf{$\Delta_{\mathrm{honest\_fusion}}$} \\ \midrule
            Honest            & Honest Candidate   & 29.6 & $+0.0$ & 49.8 & $+20.2$ & $+0.0$ \\
            Ghost             & The Ghost          & 26.7 & $-2.9$ & 34.8 & $+8.1$  & $-15.0$ \\
            Fraud             & The Fraud          & 44.3 & $+14.7$ & 47.4 & $+3.1$  & $-2.4$ \\
            Stale             & Stale Candidate    & 35.0 & $+5.4$  & 37.9 & $+2.9$  & $-11.9$ \\
            Gamer             & The Time Traveler  & 35.0 & $+5.4$  & 58.1 & $+23.1$ & $+8.3$ \\
            Squatter          & Vince Vault        & 29.6 & $+0.0$  & 29.8 & $+0.2$  & $-20.0$ \\
            Smart Squatter    & Alex Rivers Smith  & 29.6 & $+0.0$  & 49.8 & $+20.2$ & $+0.0$ \\
            Shadow            & The Shadow         & 29.6 & $+0.0$  & 30.4 & $+0.8$  & $-19.4$ \\
            Inflater          & The Inflater       & 29.6 & $+0.0$  & 46.1 & $+16.5$ & $-3.7$  \\ \bottomrule
        \end{tabular}%
    }
\end{table}





\subsubsection{Successful Behaviours}
\textbf{Squatter.} This key highlight shows the Integrity Engine successfully detecting and neutralising identity squatting. Fusion adds an insignificant $+0.2$ uplift. The $\Delta_{\mathrm{honest\_fusion}}$ reveals a $20$-point deficit compared to an identical honest candidate, proving the Integrity Engine effectively cuts off the attacker's advantage.

\textbf{Stale.} This candidate updates their CV but provides outdated repository evidence. While CV-only is unchanged, fusion adds only a minor $+2.9$ uplift. Relative to the honest candidate, they fall $-11.9$ points behind. They are not explicitly penalised, rather, Bayesian Fusion naturally reflects the lower signal strength of their time-decayed GitHub evidence.

\textbf{Ghost.} This candidate keyword-stuffs their CV while providing a public-repository that contradicts their claims. CV-only scoring partially penalises this, but their $\Delta_{\mathrm{honest\_fusion}}$ score drops heavily to $-15.0$ because Bayesian Fusion updates the Beta Distribution to reflect the contradictory GitHub signal.

\textbf{Shadow.} Unlike Ghost, this candidate does not actively contradict their CV, but their complete lack of GitHub evidence prevents them from gaining the Bayesian boost awarded to the honest candidate. While they receive a minor $+0.8$ uplift from LinkedIn, their $\Delta_{\mathrm{honest\_fusion}}$ of $-19.4$ demonstrates that missing evidence safely leaves candidates behind rather than actively penalising them.

\textbf{Inflater.} This candidate inflates their metrics by concentrating massive lines of code (LOC) into a single language, disregarding others. Because MERIT evaluates all JD requirements, neglected languages receive no demonstrated evidence. This dilutes their overall probability of competence during the Bayesian update, dragging their $\Delta_{\mathrm{honest\_fusion}}$ down to $-3.7$.

\subsubsection{Partial Success}
\textbf{Fraud.} This candidate successfully inflates their CV-only score ($+14.7$) through major false claims, while providing contradictory repository evidence. Unlike Ghost's simple keyword stuffing, Fraud overhauls their entire CV. However, Bayesian Fusion restricts their final boost to just $+3.1$ points by reflecting the low signal strength of their actual GitHub profile. Note that there is still partial success since their $\Delta_{\mathrm{honest\_fusion}}$ score is $-2.4$, Bayesian Fusion still penalises them for the false claims.

\subsubsection{Documented Failure Modes}
\textbf{Smart Squatter.} The biggest failure mode is the Smart Squatter, where the candidate bypasses the Integrity Engine by finding a highly-scored GitHub portfolio that passes the fuzzy name-match. While rare in the real world, it remains a possible attack vector.

\textbf{Gamer.} Another impractical failure is the Gamer, who manipulates public-repository metrics to trick the system into thinking they were recently active ($\Delta_{\mathrm{fusion}} = +23.1$). This highlights a critical flaw in how MERIT calculates temporal skill decay.


\subsubsection{Study 06 takeaway}
The integrity engine is able to detect and penalise most forms of realistic adversarial attacks, such as keyword stuffing, identity squatting, and repository gaming.
Note that the majority of failure modes are edge-cases that take significant effort to achieve, and are not expected to be a major concern in the real world.
For example, let us consider the Smart Squatter scenario. For a given role, let us assume the malicious candidate wishes to squat a GitHub portfolio. First,
they would need to find portfolios that would score high on MERIT, which is a somewhat feasible task if they are willing to invest the time required to do this.
From here, the actor would need to find a name that is close enough to bypass the fuzzy name-match, this is likely not a guaranteed success due to the rarity of
a candidate finding a strong portfolio that scores high on MERIT for a specific role that shares the same name as the individual. This effort would technically work, 
but in the real world, where a candidate applies to a vast array of jobs with different metrics, this would be impractical. This effort would have to be repeated for
each role family that they apply to. In our example, we consider GitHub portfolios, but this identity squatting can be applied to other data sources, such as LinkedIn as well.

\subsection{Studies 02 \& 03: Scalability \& Performance Benchmarks}
\label{subsec:studies-02-03-scalability}

Studies~02 and~03 sit outside of the research questions in Table~\ref{tab:eval-literature-traceability}. They 
assess deployment feasibility (latency and memory at batch scale) rather than whether MERIT closes the extensibility, multi-source, 
or transparency gaps from Section~\ref{sec:relevance-to-merit}.
We benchmark only against our in-repository engine variants.


The runtime and spacetime studies are performed for cohorts of size $N \in \{10, 50, 100, 200, 500\}$ against the two baselines (Traditional and Modern AI),
as well as three variants of MERIT (CV-Only, Full without explainability, Full with explainability). In Study~02, we initialise each
engine once and then measure the wall-clock time for each engine under each cohort size. For study~03, we separate the initialisation 
cost from the incremental batch overhead per candidate. We run these benchmarks on Windows~11 (AMD Ryzen~7 7800X3D, Python~3.11.9).
Scoring is a purely CPU-bound operation, we do not leverage GPU acceleration or other hardware accelerations. We also do not
account for costs such as PDF ingestion, API calls and network latency from the interactions with our Supabase database.
The results are presented in Table~\ref{tab:runtime-results} and~\ref{tab:study03-memory}, 
for a graphical representation, see Appendix~\ref{appendix:fairness-scalability-charts}.

\begin{table}[htbp]
    \centering
    \caption[Study~02: Batch ranking latency]{Study~02: batch ranking latency (seconds) vs.\ cohort size $N$.}
    \label{tab:runtime-results}
    \vspace{1.5ex}
    \small
    \setlength{\tabcolsep}{6pt}
    \begin{tabular}{r c c c c c}
        \toprule
        \textbf{$N$} & \textbf{Trad. ATS} & \textbf{Mod. AI ATS} & \textbf{MERIT CV} & \textbf{MERIT Full} & \textbf{MERIT Expl.} \\ \midrule
        10  & 0.003  & 4.33   & 0.090 & 0.092 & 0.513 \\
        50  & 0.012  & 21.14  & 0.435 & 0.465 & 2.588 \\
        100 & 0.026  & 42.18  & 0.890 & 0.963 & 5.169 \\
        200 & 0.040  & 85.03  & 1.789 & 1.862 & 10.48 \\
        500 & 0.088  & 212.62 & 4.409 & 4.842 & 26.06 \\ \bottomrule
    \end{tabular}
\end{table}

\begin{table}[htbp]
    \centering
    \caption[Study~03: Engine memory]{Study~03: engine memory (MB). \textit{init.}\,=\,initialisation (static load at startup); remaining rows ($N$): incremental batch overhead after \texttt{gc.collect()} (not cumulative with static load).}
    \label{tab:study03-memory}
    \vspace{1.5ex}
    \small
    \setlength{\tabcolsep}{6pt}
    \begin{tabular}{r c c c c c}
        \toprule
        \textbf{$N$} & \textbf{Trad. ATS} & \textbf{Mod. AI ATS} & \textbf{MERIT CV} & \textbf{MERIT Full} & \textbf{MERIT Expl.} \\ \midrule
        \textit{init.} & 0.10 & 398.0 & 367.1 & 367.5 & 367.7 \\ \addlinespace
        10  & 0.02 & 113.0 & 54.4 & 54.5 & 53.8 \\
        50  & 0.03 & 116.7 & 56.3 & 56.7 & 56.7 \\
        100 & 0.05 & 118.5 & 58.8 & 59.1 & 59.4 \\
        200 & 0.06 & 119.7 & 63.3 & 63.5 & 65.0 \\
        500 & 0.12 & 119.1 & 73.9 & 74.6 & 76.5 \\ \bottomrule
    \end{tabular}
\end{table}

At $N = 500$, we see that MERIT Full is able to score the batch in under $4.4$ seconds, compared to $212.6$ seconds 
for the Modern AI ATS. Note the introduction of explainability, specifically Shapley values results in a runtime of 
$26.1$ seconds, leading to a $5.93$x slowdown compared to MERIT without explainability. There are $(2^3 = 8)$ possible 
coalitions of sources, and this will scale exponentially with the number of sources, not cohort size. Multi-source
fusion adds negligible latency and RSS overhead compared to CV only ($+0.43$ seconds and under $1$ MB at $N = 500$).
The majority of static load can be attributed to the SBERT initialisation (398.0 MB for Modern AI ATS, 367.1 MB for MERIT).
MERIT scales linearly with $N$ ($54.4$ MB to $73.9$ MB) because the harness retains per-candidate payloads, however,
this is still relatively modest for desktop batch screening.
